\section{Introduction}

\paperfig{fig_overview.png}{F1}{\textbf{Overview.} Top: TeachQuiz asks the student to forget the topic; instead it refuses, it already knows
MMMC, and the grader lost answers. Our measurement makes the student choose, measures what it knew, and uses invented
concepts. Bottom: how we repair a video with learner feedback, and how we evaluate the repair.}{overview}

Code-generated teaching videos are made in three steps: a planner writes a lesson, a coder writes Manim~\cite{manim}
scenes, and a renderer turns them into video. Such videos are now judged by whether a model \emph{learns} from them.
Code2Video's TeachQuiz~\cite{code2video} tells a vision--language student to forget a topic. It then measures the
student's quiz score without the video ($S_1$) and with it ($S_2$), and reports $S_2 - S_1$.

We set out to use this measure for a stronger claim: that a generated video can be \emph{repaired} where a simulated
learner failed. The repair traces each failure to the section that should have taught it and rewrites only that
section. The test is whether teaching improves on unseen questions written by other people.

A pilot on 12 MMMC topics found no signal (Table~\ref{tab:pilot}). The learner repair barely moved TeachQuiz, and
\PilotTies{} of 12 topics were ties with the base video. When we looked for the cause, we found that the measure was
at fault, not only the method. It failed in three ways that add up:
\begin{enumerate}
  \item \textbf{The grader loses answers.} The questions are named \texttt{t1..t5}, but the student copies the
  example id from the answer template (\texttt{q1}). Matched by id, all its answers are lost. This happened in \OrderRuns{} of \AllRuns{} student
  runs (\OrderShare\%) on TEST, and not on DEV, whose ids match the template.
  \item \textbf{The student refuses instead of forgetting.} Told to forget, it
  answers ``I don't know'' instead of guessing (\emph{refusal}). Without the video it answers NULL on \NullNoVideo\% of
  items; with the video, still on \NullWithVideo\%. So $S_1 = 0$ measures obedience, and $S_2$ measures when the student
  feels allowed to answer.
  \item \textbf{The student already knows the answers.} Made to choose and shown nothing, it answers
  \mmmcpriorGemmaAcc\% of all MMMC TEST items correctly. Its mean confidence in the right letter is
  \mmmcpriorGemmaP\%, and it gets \mmmcpriorGemmaAllRight{} of \mmmcpriorGemmaTopics{} topics entirely right. A 7B
  student does the same (\mmmcpriorQwenAcc\%). We call this a \emph{prior ceiling}: no video can show a gain on items
  the student already knows.
\end{enumerate}

\paragraph{Contributions.}
\begin{itemize}
  \item An audit of TeachQuiz that separates the three failure modes above, with the evidence for each.
  \item D-VGAIN, a replacement measure. It makes the student choose, subtracts what the student knew before, and
  compares the moving video with its start and end frames. A \emph{metric ladder} adds these parts one at a time, giving ten
  versions of the measure (M0 to M9). Each version must pass four tests: it gives a score, it separates a topic's own video from another topic's, it repeats, and it
  breaks ties (\S\ref{sec:ladder}).
  \item NOVEL-12, a benchmark of invented concepts that no model has seen, so a gain cannot come from memory
  (\S\ref{sec:novel}).
  \item An admission check for videos and two model-free checks: one for when objects are on screen, one for what
  they are attached to. With them we show that the layout score misses faults a viewer sees (\S\ref{sec:quality}).
  \item A pipeline ablation that adds five repair features one at a time: an enforced gate, code checks before
  rendering, finding the weak section like a debugger, a leak guard, and minimal edits. It is built and its rules are
  fixed in advance; its evaluation is future work (\S\ref{sec:ablation}).
\end{itemize}
